% Fazit und Ausblick — claude-tunnel bachelor thesis.
\chapter{Fazit und Ausblick}
\label{ch:fazit}

\section{Zusammenfassung}
\label{sec:fazit-zusammenfassung}

Diese Arbeit hat einen providerblinden, zensurresistenten Netzwerk-Tunnel
entworfen, als Rust-Prototyp aus fuenf Crates implementiert und in einem rein
containerbasierten Testbett evaluiert. Der Betreiber der Vermittlungsebene relayt
ausschliesslich Chiffretext: die Nutzlast ist mit dem Noise-Protokoll ueber
\gls{ac:quic} Ende-zu-Ende verschluesselt und an die gepinnte Origin-Identitaet
gebunden, die Adressierung erfolgt ueber opake Tokens statt Hostnamen. Ein
Proof-of-Work-Rendezvous daempft Missbrauch, ein direkter Peer-to-Peer-Pfad
umgeht bei Erreichbarkeit den Relay, und ein
\gls{ac:tls}-ueber-\gls{ac:tcp}-Fallback traegt denselben Tunnel bei blockiertem
\gls{ac:udp}.

Auf die Forschungsfragen (Kapitel~\ref{ch:einleitung}) antwortet die Arbeit wie
folgt: Ein solcher Tunnel ist realisierbar (FF1) -- Architektur und
Kernmechanismen wurden bis hin zum verschluesselten Roundtrip durch echte
Komponenten belegt, einschliesslich der guthabengedeckten Token-Ausgabe. Der
Latenz-Overhead ist gering und vorhersagbar (FF2): der Eigenanteil liegt bei rund
8\,ms, und der Roundtrip skaliert linear mit der Netzverzoegerung
($\text{RT} \approx 8{,}8 + 6{,}1 \cdot d$). Paketverlust trifft primaer den Tail,
und die Betriebsart ist bei verlustfreiem Netz nicht latenzbestimmend (FF3).

Der wesentliche Beitrag ist die \emph{Kombination} von struktureller
Providerblindheit, Ende-zu-Ende-Pinning, Proof-of-Work-Rendezvous, direktem Pfad
und Transport-Fallback in einem einzigen, lauffaehigen System -- ergaenzt um eine
reproduzierbare, statistisch charakterisierte Vermessung des dadurch entstehenden
Overheads.

\section{Ausblick}
\label{sec:fazit-ausblick}

Aus den in Kapitel~\ref{ch:diskussion} benannten Grenzen ergeben sich mehrere
Anschlussrichtungen:

\begin{itemize}
  \item \textbf{Traffic-Analyse-Resistenz.} Padding und Ratenglaettung koennten
    Groessen- und Timing-Seitenkanaele gegen einen beobachtenden Vermittler
    verdecken -- als orthogonale Schicht mit eigenem, zu vermessendem Overhead.
  \item \textbf{Sybil-Oekonomie.} Ueber den reinen Token-Preis hinausgehende
    Kosten pro schaedlicher Aktion oder pseudonyme Reputation koennten den
    finanzierten Angreifer (A4) adressieren, ohne die Pseudonymitaet aufzugeben.
  \item \textbf{NAT-Traversal unter realen Bedingungen.} Ein vollstaendiges
    \gls{ac:ice}-Verfahren und ein Testbett mit emuliertem symmetrischem
    \gls{ac:nat} wuerden das Hole-Punching ueber die flache Bridge hinaus
    belastbar machen.
  \item \textbf{Overhead-Senkung.} Verbindungs-Wiederverwendung (0-RTT,
    Keep-Alive) koennte den durch den Frisch-Dial dominierten Handshake-Anteil
    deutlich reduzieren.
  \item \textbf{Breitere Parameterstudie.} Groesseres $n$, eine
    Bandbreiten-Achse und eine ausgefahrene PoW-Schwierigkeitsachse wuerden die
    Tail-Aussagen und einen etwaigen Modus-Effekt unter Verlust schaerfen.
\end{itemize}

Insgesamt zeigt die Arbeit, dass providerblinde Vermittlung nicht als
organisatorisches Versprechen, sondern als strukturelle Eigenschaft eines
praktisch betreibbaren Tunnels realisierbar ist -- zu einem geringen und
vorhersagbaren Latenzpreis.
